% ============================================================
\chapter*{Preface}
\addcontentsline{toc}{chapter}{Preface}
% ============================================================

\section*{Why 90\,\% of ML Projects Fail in Deployment}

A recurring finding across industry is alarming: according to Gartner (2022),
only 54\,\% of machine learning (ML) models make it from prototype to
production, and among those that do, many are retired within the first few
months. VentureBeat reported in 2019 that 87\,\% of data science projects
never reach production. This is not a modelling problem---it is an
\textbf{engineering} problem.

\begin{center}
\begin{tikzpicture}[
  block/.style={rectangle, draw, rounded corners, minimum width=2.8cm,
    minimum height=1cm, align=center, font=\small},
  arrow/.style={-{Stealth[length=3mm]}, thick},
  fail/.style={-{Stealth[length=3mm]}, thick, red!70!black, dashed}
]
  \node[block, fill=blue!10] (data) {Data};
  \node[block, fill=blue!10, right=1.2cm of data] (notebook) {Notebook\\Experimentation};
  \node[block, fill=green!10, right=1.2cm of notebook] (model) {Trained\\Model};
  \node[block, fill=orange!10, right=1.2cm of model] (deploy) {Deployment\\Production};
  \node[block, fill=red!10, right=1.2cm of deploy] (monitor) {Monitoring\\Maintenance};

  \draw[arrow] (data) -- (notebook);
  \draw[arrow] (notebook) -- (model);
  \draw[fail] (model) -- node[above, font=\scriptsize, red!70!black] {Frequent failure} (deploy);
  \draw[arrow] (deploy) -- (monitor);

  % Valley of death
  \node[below=0.6cm of model, font=\small\itshape, text=red!70!black] {``Valley of Death''};
\end{tikzpicture}
\end{center}

The main causes of these failures are:
\begin{itemize}
  \item \textbf{Technical debt}\index{technical debt}: ML code accumulates
    hidden complexity (data dependencies, feedback loops, unversioned
    configurations)~\cite{sculley2015}.
  \item \textbf{Irreproducibility}: impossible to recreate a colleague's
    results, even with the same code, because the environment, data, or
    hyperparameters differ.
  \item \textbf{Organizational gap}: data scientists work in isolated
    notebooks, ML engineers do not understand modelling choices, and ops
    teams do not know how to monitor a model.
  \item \textbf{Lack of monitoring}: a deployed model without surveillance
    silently drifts (data drift, concept drift).
  \item \textbf{No reproducible pipeline}: every step (preprocessing,
    training, evaluation) is executed manually.
\end{itemize}

\section*{What MLOps Solves}

\textbf{MLOps}\index{MLOps} (Machine Learning Operations) is a set of
practices that unify ML development (Dev) and operations (Ops) to automate,
standardize, and make reliable the lifecycle of models. It is the equivalent
of DevOps for machine learning, with specific challenges: data versioning,
experiment tracking, result reproducibility, drift monitoring.

\begin{center}
\begin{tikzpicture}[
  phase/.style={rectangle, draw, rounded corners, minimum width=2.2cm,
    minimum height=0.9cm, align=center, font=\small\bfseries},
  arrow/.style={-{Stealth[length=2.5mm]}, thick, steelblue}
]
  \node[phase, fill=blue!15] (design) {Design};
  \node[phase, fill=blue!15, right=0.8cm of design] (dev) {Development};
  \node[phase, fill=green!15, right=0.8cm of dev] (train) {Training};
  \node[phase, fill=orange!15, right=0.8cm of train] (deploy) {Deployment};
  \node[phase, fill=red!15, right=0.8cm of deploy] (monitor) {Monitoring};

  \draw[arrow] (design) -- (dev);
  \draw[arrow] (dev) -- (train);
  \draw[arrow] (train) -- (deploy);
  \draw[arrow] (deploy) -- (monitor);
  \draw[arrow, bend right=30] (monitor) to node[below, font=\scriptsize] {re-training} (train);
\end{tikzpicture}
\end{center}

\section*{Course Organization}

This course is designed for Master-level students in data science, machine
learning, or applied mathematics. It assumes a solid foundation in Python,
familiarity with the Linux/Unix command line, and knowledge of machine
learning.

Each chapter follows a consistent pedagogical structure:
\begin{enumerate}
  \item \textbf{Motivation and concepts}: why this tool or practice is
    necessary.
  \item \textbf{Hands-on tutorial}: working code that the student can run
    on their own machine.
  \item \textbf{Best practices and anti-patterns}: what to do and what to
    avoid.
  \item \textbf{Tool comparison}: comparison tables when multiple solutions
    exist.
  \item \textbf{Mini-project}: integrative exercise using tools from
    previous chapters.
  \item \textbf{Exercises}: graded by difficulty
    ($\bigstar$ setup/config, $\bigstar\bigstar$ pipeline building,
    $\bigstar\bigstar\bigstar$ end-to-end project).
  \item \textbf{Cheatsheet}: summary of essential commands.
\end{enumerate}

\vfill
\noindent\textit{Happy reading and happy deploying!}
